% =====================================================================
% EdgeFlow 用户手册 · 第 5 章 边缘节点与自治运行
% 事实依据：docs/KEADM.md、docs/RELEASE-NOTES-v0.1.0.md、edge/pkg/edgehub、
%           edge/pkg/edged、cloud/pkg/nodecontroller、cmd/edgecore
% =====================================================================
\chapter{边缘节点与自治运行}
\label{ch:5}

EdgeFlow 的边缘侧由一个常驻进程 \textbf{edgecore}（边缘核心）承载。它运行在
边缘节点上，负责三件事：与云端建立并维持管理通道、按云端期望状态调谐本机
容器、以及在本机维护设备影子与元数据。本章说明 edgecore 的接入、通信、
自治行为，以及如何配置它。

\section{接入与连接管理}

\subsection{注册}

edgecore 启动后第一条消息是 \textbf{Register}（注册），云端 CloudHub 校验后
回 \textbf{RegisterAck}（注册确认）。注册通过后，节点在云端注册表中可见，
状态为 \texttt{Ready}。节点 ID 由环境变量 \texttt{EDGEFLOW\_EDGECORE\_NODE\_ID}
或配置文件（\texttt{config/edgecore.json} 的 \texttt{nodeID} 字段）指定，环境变量
优先；均未设置时使用默认约定 \texttt{edge-\allowbreak<主机名>}。

\begin{lstlisting}
# 指定节点 ID 与云端地址启动 edgecore
export EDGEFLOW_EDGECORE_NODE_ID=edge-worker-01
export EDGEFLOW_EDGECORE_CLOUD_ADDR=ws://192.168.1.10:10000/v1/edge
bin/edgecore
\end{lstlisting}

\subsection{心跳保活}

注册完成后，edgecore 以 \textbf{30 秒}为周期向云端发送心跳（Heartbeat）。
云端 NodeController 每 \textbf{30 秒}扫描一次注册表，若某节点心跳停滞超过
\textbf{180 秒}（约 6 个心跳周期），将该节点标记为 \texttt{Offline}。
因此：

\begin{itemize}
  \item 心跳周期：固定 \texttt{30s}（云边契约常量，不可配置）；\texttt{EDGEFLOW\_EDGECORE\_REPORT\_INTERVAL} 仅控制 Pod 状态上报周期；
  \item 云端超时阈值：\texttt{EDGEFLOW\_CLOUDCORE\_NODE\_TIMEOUT}，默认 \texttt{180s}；
  \item 云端扫描周期：\texttt{EDGEFLOW\_CLOUDCORE\_NODE\_SCAN\_INTERVAL}，默认 \texttt{30s}。
\end{itemize}

\subsection{通道压缩}

云边通道的 gzip 压缩\textbf{默认开启}：edgecore 在 Register 消息中声明压缩
能力（\texttt{compression: "gzip"}），云端在 RegisterAck 中回带确认，协商
成功后双向启用。\texttt{config/cloudcore.json} 中设置 \texttt{compress:false}
可关闭：云端不再回带压缩能力，边缘自动回落明文传输，与旧版本互操作兼容。

\subsection{断线重连}

网络抖动或云端重启导致连接断开时，edgecore \textbf{自动重连}，无需人工干预：

\begin{itemize}
  \item 指数退避：\texttt{1s/2s/4s/…}，上限 \texttt{60s}；
  \item 重连成功后重新执行注册流程，退避计时重置；
  \item 重连期间本地功能不受影响（见第~\ref{ch:5}~章第~\ref{sec:5-autonomy}~节）。
\end{itemize}

\section{可靠投递}

云边之间的下行消息（Pod 配置、设备指令等）使用\textbf{可靠投递}语义，保证
云端能确认消息确实被边缘处理：

\begin{enumerate}
  \item 云端发送消息并登记到 \textbf{pending}（待确认）队列；
  \item 边缘收到并处理成功后回 \textbf{Ack}；
  \item 云端收到 Ack 后移出 pending 队列；
  \item 超时（默认单次 5 秒）未确认，使用\textbf{同一消息 ID} 重试，最多
        3 次；仍失败则向 API 调用方返回 504。
\end{enumerate}

边缘侧对重复消息做\textbf{幂等去重}：同一消息 ID 处理过一次后不再重复执行，
因此重试不会造成重复下发。API 层表现见第~\ref{ch:9}~章与附录~B。

\section{Edged 应用调谐}
\label{sec:5-edged}

\textbf{Edged} 是 edgecore 内置的容器运行时管理模块，对标 KubeEdge 的 Edged。
它采用\textbf{声明式调谐}模型：云端下发期望状态（Pod 定义）后，Edged 周期性
比对期望与实际情况，差异即收敛。

\begin{itemize}
  \item 调谐周期：\texttt{EDGEFLOW\_EDGECORE\_RECONCILE\_INTERVAL}，默认 \texttt{5s}；
  \item 启动时立即调谐一次，之后按周期循环；
  \item 调谐数据来自 MetaManager 落盘数据（SQLite），因此断网时调谐照常进行。
\end{itemize}

\subsection{多副本 Pod}

云端下发的 Pod 支持 \texttt{replicas} 多副本。Edged 每轮调谐将本机实际运行的
副本数收敛到期望值：

\begin{itemize}
  \item 实际副本 \textbf{少于}期望 → 逐个补齐拉起；
  \item 实际副本 \textbf{多于}期望 → 逐个清理收缩；
  \item 每轮调谐只处理一个副本的差异（批大小 1），逐轮滚动，避免瞬间抖动。
\end{itemize}

\subsection{健康自愈与 CrashLoopBackOff}

Edged 每轮调谐都会检查容器运行状态（\texttt{Inspect}）。容器非 Running
（异常退出、被杀死等）时自动重启，无需人工介入：

\begin{itemize}
  \item 每次重启累加 \texttt{RestartCount} 并上报云端；
  \item 连续异常重启达到\textbf{3 次}阈值后进入 \textbf{CrashLoopBackOff}
        退避：退避时长 \textbf{30 秒}，退避期间不再反复拉起；
  \item 容器稳定运行 \textbf{60 秒}后重置计数，恢复正常自愈节奏。
\end{itemize}

典型表现：退出型镜像（如 \texttt{busybox} 执行完即退出）会被反复拉起并进入
退避，属预期行为；业务容器持续崩溃时请查看容器日志定位根因（见
第~\ref{ch:9}~章）。

\subsection{镜像漂移检测}

Edged 在调谐时会比对\textbf{期望镜像}与\textbf{实际运行镜像}：

\begin{itemize}
  \item 已运行的容器镜像与期望不一致 → 判定为\textbf{镜像漂移}，自动重建：
        停止旧容器 → 用期望镜像重新拉起；
  \item 同一轮只重建 1 个漂移容器，逐轮滚动完成全部收敛；
  \item 该机制同时是滚动更新的基础：更新 Pod 定义中的镜像后，Edged
        自动逐副本完成替换。
\end{itemize}

\section{断网自治}
\label{sec:5-autonomy}

\textbf{自治}是边缘计算的核心价值：云端不可达时，边缘业务不能停。

\begin{enumerate}
  \item 云端断开后，已下发的容器\textbf{持续运行}，调谐、设备采集、本地
        指令照常工作；
  \item 断网语义验证：自动化用例以\textbf{60 秒短时断网}模拟验证；
        \textbf{30 分钟}级长时运行语义（M2 验收口径）需在真实环境长跑确认；
  \item 重连成功后，edgecore 重新注册，并将离线期间的 Pod 状态、
        设备上报\textbf{自动同步}到云端，云端视图恢复一致。
\end{enumerate}

注意：云端状态为\textbf{内存态}，云端进程重启后节点需重新注册（见
第~\ref{ch:9}~章）。边缘侧元数据由 MetaManager 持久化到本机 SQLite，断网
期间不丢失。

\section{edgecore 配置（环境变量 + 配置文件）}

edgecore 采用\textbf{环境变量 + 配置文件}双轨配置，加载优先级从高到低为：
\textbf{环境变量 > 配置文件 > 默认值}。配置文件为 JSON 格式，默认路径
\texttt{config/edgecore.json}（可用 \texttt{--config} 标志覆盖），所有字段
可选（缺省回落默认值），支持 \texttt{cloudAddr}、\texttt{nodeID}、
\texttt{podReportInterval}、\texttt{deviceReportInterval} 与
\texttt{reconcileInterval} 五个字段，例如：

\begin{lstlisting}
{
  "cloudAddr": "ws://192.168.1.10:10000/v1/edge",
  "nodeID": "edge-worker-01",
  "podReportInterval": "30s",
  "deviceReportInterval": "30s",
  "reconcileInterval": "5s"
}
\end{lstlisting}

环境变量全集见表~\ref{tab:5-env}；时长类变量（如调谐/上报周期）合法范围为
\textbf{1s 至 10m}，超出范围或非法值回退默认值并打印告警。\textbf{敏感配置
（接入令牌 \texttt{EDGEFLOW\_EDGECORE\_TOKEN} 等）仍仅支持环境变量，不入
配置文件}，防止令牌随文件分发泄露。

cloudcore 采用同构机制：配置文件 \texttt{config/cloudcore.json}（字段
\texttt{port}、\texttt{hubPort}、\texttt{compress}，全部可选）+ \texttt{--config}
覆盖；优先级为命令行 > 环境变量 > 配置文件 > 默认值。

\begin{table}[htbp]
\centering
\caption{edgecore 环境变量}
\label{tab:5-env}
\begin{tabularx}{\textwidth}{>{\raggedright\arraybackslash}p{4.5cm} >{\raggedright\arraybackslash}p{2.8cm} X}
\toprule
\textbf{变量} & \textbf{默认值} & \textbf{说明} \\
\midrule
\texttt{EDGEFLOW\_EDGECORE\_NODE\_ID} & \texttt{edge-\allowbreak<主机名>} & 节点 ID，注册用 \\
\texttt{EDGEFLOW\_EDGECORE\_CLOUD\_ADDR} & \texttt{ws://127.0.0.1:10000/v1/edge} & 云端 CloudHub 地址 \\
\texttt{EDGEFLOW\_EDGECORE\_TLS} & 关 & \texttt{on} 时启用 mTLS（地址自动 \texttt{wss://}） \\
\texttt{EDGEFLOW\_EDGECORE\_CERT\_DIR} & \texttt{data/certs/} & 证书目录 \\
\texttt{EDGEFLOW\_EDGECORE\_DB\_PATH} & 内置默认 & MetaManager SQLite 路径 \\
\texttt{EDGEFLOW\_EDGECORE\_MQTT\_ADDR} & \texttt{tcp://127.0.0.1:1883} & 本机 MQTT broker 地址 \\
\texttt{EDGEFLOW\_EDGECORE\_RECONCILE\_INTERVAL} & \texttt{5s} & Edged 调谐周期 \\
\texttt{EDGEFLOW\_EDGECORE\_REPORT\_INTERVAL} & \texttt{30s} & Pod 状态上报周期 \\
\texttt{EDGEFLOW\_EDGECORE\_DEVICE\_REPORT\_INTERVAL} & \texttt{30s} & 设备状态上报周期 \\
\texttt{EDGEFLOW\_EDGECORE\_MQTT\_CONNECT\_TIMEOUT} & \texttt{5s} & 单次 MQTT 建连超时 \\
\texttt{EDGEFLOW\_EDGECORE\_TOKEN} & 空 & 接入令牌（keadm join 写入；仅环境变量，不入配置文件） \\
\bottomrule
\end{tabularx}
\end{table}

生产环境建议使用 \texttt{keadm join} 生成 \texttt{edgecore.env} 与
\texttt{edgecore.service}（systemd 单元），由安装脚本统一部署，避免手工
拼写错误（见第~\ref{ch:7}~章升级回滚时的产物管理）。\texttt{keadm join}
产物与配置文件\textbf{双轨并存}：两者同时设置同一项时，以环境变量为准。

\section{配置热重载}

cloudcore 与 edgecore 均支持\textbf{配置热重载}，无需重启进程即可让部分
配置生效：

\begin{itemize}
  \item \textbf{触发方式}：向进程发送 \textbf{SIGHUP} 信号立即强制重载一次
        （不检查变更）；此外进程每 \textbf{60 秒}检查一次配置文件，
        mtime/size 变化即触发重载；多触发源串行化，并发安全；
  \item \textbf{重载失败保持旧配置}（fail-safe）：配置文件解析失败、文件
        缺失或端口绑定失败时，本次重载被整体拒绝，进程按旧配置继续运行。
\end{itemize}

不同配置项的热生效边界见表~\ref{tab:5-reload}：

\begin{table}[htbp]
\centering
\caption{配置热重载生效边界}
\label{tab:5-reload}
\begin{tabularx}{\textwidth}{l l X}
\toprule
\textbf{组件} & \textbf{配置项} & \textbf{热重载行为} \\
\midrule
cloudcore & \texttt{port}（HTTP） & \textbf{热切换}：先绑定新端口，成功后关闭旧监听；新端口绑定失败则拒绝本次重载，保持旧配置 \\
cloudcore & \texttt{hubPort} / \texttt{compress} & \textbf{需重启}：热重载时打印告警并保持旧值（CloudHub 不支持运行期重建监听；压缩在连接注册时协商） \\
edgecore & \texttt{podReportInterval} / \texttt{deviceReportInterval} & \textbf{热生效}：重载后按新周期上报，无需额外动作 \\
edgecore & \texttt{cloudAddr} / \texttt{nodeID} / \texttt{reconcileInterval} & \textbf{需重启}：热重载时打印告警并保持旧值 \\
\bottomrule
\end{tabularx}
\end{table}

生产建议：修改\textbf{需重启}类配置项后，在维护窗口重启对应进程使其生效；
热生效类配置可即时下发，适合运行期调整上报频率。
